\section{Introduction}
\label{sec:intro}

Automated code review is moving from passive summarisation toward agents that
fetch artifacts, run or read tests, and emit binding accept/reject decisions.
Evaluating such agents is subtle because a correct-looking label can arise from
insufficient or corrupted evidence, and a withheld decision---asking for more
evidence---can be the \emph{correct} response when the granted evidence is
non-decisive. A benchmark that scores only \ACCEPT/\REJECT\ against a single
oracle cannot tell these apart and will systematically reward
over-commitment.

We organise the evaluation around one question:
\begin{quote}
\emph{Can code-review agents distinguish evidence integrity from evidence
sufficiency?}
\end{quote}

\emph{Evidence integrity} asks whether the artifacts an agent sees are
authentic, fresh, linked to the patch, and complete. \emph{Evidence
sufficiency} asks whether those artifacts, even if perfectly intact, are
decisive for the review verdict. These are orthogonal: a restricted evidence
view can be 100\,\% authentic yet genuinely insufficient (the correct output is
\NEED), while a sufficient view can be misread (a visible failure can be
ignored).

GroundedReview makes four contributions.
\begin{enumerate}[leftmargin=1.6em,itemsep=2pt,topsep=3pt]
  \item \textbf{Benchmark formulation.} A three-way review task
  (\ACCEPT/\REJECT/\NEED) with an explicit distinction between
  \emph{Oracle Gold} and \emph{Condition Gold}, and a \emph{scoped} notion of
  acceptability tied to declared tests rather than whole-project correctness.
  \item \textbf{Evidence-controlled evaluation.} Five evidence conditions per
  issue that vary (i) static versus active tool access and (ii) restricted
  versus decisive test outcomes, with machine-checked non-interference so that
  restricted surfaces cannot leak the verdict.
  \item \textbf{Empirical analysis.} A corrective replication that isolates a
  protocol-construct artifact in our own earlier instrument, evidence of
  unsupported commitment under restricted evidence, and a traceable case of a
  sufficient-but-ignored failure---all reported with explicit denominators,
  issue-level units, and confidence bounds, alongside an independently
  constructed validation set, since reviewed by an independent human
  annotator.
  \item \textbf{Controlled causal evidence (G4).} A pre-registered,
  within-issue $2\times2$ intervention over 18 frozen pairs using a yoked
  fixed evidence transcript that holds linkage constant and independently
  manipulates decisive-outcome visibility and integrity attestation, analysed
  with paired exact tests and independent contexts, which isolates a
  model-specific causal effect of attestation on unsupported acceptance.
\end{enumerate}

We emphasise what the present evidence does \emph{not} support:
population-level false-acceptance rates (our oracle-\REJECT\ sample is three,
later eight with only two rejects, and 18 with five rejects in G4), causal
attribution of the cross-version change to a single fix (the G4 intervention
separately supplies the one controlled, within-issue causal estimate in the
paper), or any claim that the scoped \ACCEPT\ issues passed the entire SymPy
test suite.
